\subsection{Positive elements}

Let \(G\) be a pre-ordered group with pre-order relation \(\preccurlyeq\) ; if \(0 \preccurlyeq x\) and \(0 \preccurlyeq y\) then we deduce that \(y \preccurlyeq x + y\) by (POM), and so \(0 \preccurlyeq x + y\) by transitivity; this says that the set \(G_{+}\) of elements \(x \in G\) such that \(0 \preccurlyeq x\) is closed under addition; moreover, the relation \(x \preccurlyeq y\) is equivalent to \(0 \preccurlyeq y - x\) , that is to \(y - x \in G_{+}\) . Conversely:

PROPOSITION 3. --- If P is a subset of an abelian group G, containing 0 and such that \(P + P \subset P\) , then the relation \(y - x \in P\) is a pre-order relation compatible with the group structure of G. For this relation to make G an ordered group it is necessary and sufficient to have \(P \cap (-P) = \{0\}\) ; for G to be a totally ordered group under this ordering it is necessary and sufficient to have in addition \(P \cup (-P) = G\) .

It is immediate that the relation \(y - x \in P\) is reflexive and transitive, and (if it is written \(x \preccurlyeq y\) ) satisfies the axiom (POM). To prove the second assertion it is enough to remark that \(P \cap (-P)\) is the subgroup \(G'\) consisting of elements \(x\) such that \(0 \preccurlyeq x\) and \(x \preccurlyeq 0\) . Finally, to say that \(G\) is totally ordered means that, for each pair of elements \(x, y\) of \(G\) , at least one of the elements \(x - y, y - x\) belongs to \(P\) , which completes the proof.

DEFINITION 3. --- In an ordered group G, any element x such that \(0 \leqslant x\) (resp. \(x \leqslant 0\) ) is called a positive (resp. negative) element.

Note that 0 is the only element which is both positive and negative; any element x such that 0 \textless{} x (resp. x \textless{} 0) is called strictly positive (resp. strictly negative).

Example. --- In the additive group \(Z \times Z\) , let P be the set of elements \((x, y)\) satisfying two inequalities \(ax + by \geqslant 0\) and \(cx + dy \geqslant 0\) , where a, b, c, d are integers (* or real numbers *) such that \(ad - bc \neq 0\) ; the « cone » P satisfies the first two conditions of Prop. 3. In this way various orderings compatible with the group structure of \(Z \times Z\) can be defined; the group is not totally ordered under any of these orderings.

Remark. --- By virtue of the condition \(P + P \subset P\) , the relation \(x \geqslant 0\) in an ordered group implies that \(nx \geqslant 0\) for every natural number n.~If in addition the positive element x of G has finite order n, then \(-x = (n - 1)x\) is positive; since

\[
\mathbf {P} \cap (- \mathbf {P}) = \{0 \},
\]

this implies that \(x = 0\) . In particular, if every element of \(\mathbf{G}\) has finite order, then \(\mathbf{P} = \{0\}\) ; the relation \(x \leqslant y\) is then equivalent to \(x = y\) (the discrete ordering).

